\chapter{正規表現の基本}
\label{app:regex}

正規表現は、文字列の集合を短く記述するための書き方である。\code{grep}、\code{awk}、Python の \code{re} モジュールなどで共通して使われ、ログの抽出、ファイル名の検査、検出器 ID の選別、時刻文字列の確認といった場面で役に立つ。ここでは、本文で触れた \code{grep} や \code{awk} の延長として、解析で本当に必要になる範囲をまとめる。

\section{正規表現とは何か}

\subsection{文字どおりの一致とメタ文字}

最初の一歩は、普通の文字列検索と正規表現の違いを区別することである。\code{det\_A} という文字列だけを探したいなら、文字をそのまま並べればよい。

\begin{lstlisting}[numbers=none, xleftmargin=2\zw]
$ grep "det_A" events.txt
\end{lstlisting}

これに対して、\code{.} や \code{*} のような特別な意味を持つ記号を使うと、「どの文字でもよい」「0 回以上くり返す」といった条件を書けるようになる。こうした特別な意味を持つ文字をメタ文字という。

\begin{lstlisting}[numbers=none, xleftmargin=2\zw]
.
*
+
?
^
$
[]
()
|
\end{lstlisting}

たとえば \code{det\_.} は \code{det\_A} にも \code{det\_B} にも一致する。\code{.} は「任意の 1 文字」を表すからである。逆に、ピリオドそのものを探したいなら、前にバックスラッシュを付けて \texttt{\textbackslash.} のようにエスケープしなければならない。これは \code{+} や \code{?} でも同じである。

\subsection{ワイルドカードとの違い}

シェルのワイルドカードと正規表現は似て見えるが、同じものではない。\code{ls *.txt} の \code{*} は「任意の文字列」を表すが、正規表現の \code{*} は「直前の要素の 0 回以上のくり返し」を表す。したがって、正規表現で「任意の文字列」を書きたいときは \code{.*} と書く。

\begin{table}[hbtp]
\centering
\caption{ワイルドカードと正規表現の違い}
\label{tab:wildcard-vs-regex}
\begin{tabular}{lll}
\toprule
場面 & 記法 & 意味 \\
\midrule
シェル & \code{*.txt} & 任意の名前で \code{.txt} で終わるファイル \\
正規表現 & \texttt{.*\textbackslash.txt\$} & 任意の文字列で始まり \code{.txt} で終わる文字列 \\
\bottomrule
\end{tabular}
\end{table}

表~\ref{tab:wildcard-vs-regex} は、見た目が似ていて混同しやすい 2 つの記法を並べたものである。この違いを混同すると、\code{grep "*.txt"} のように書いて期待通りに動かない。シェルが展開する記法なのか、正規表現エンジンが読む記法なのかを意識することが重要である。

\section{よく使う部品}

\subsection{アンカー（行頭と行末）}

\code{\textasciicircum} は行頭、\code{\$} は行末を表す。したがって、文字列の一部ではなく「行全体がその形である」ことを確かめたいときに使う。

\begin{lstlisting}[numbers=none, xleftmargin=2\zw]
$ grep '^ERROR' run.log
$ grep 'txt$' filelist.txt
$ grep '^det_A$' detectors.txt
\end{lstlisting}

1 行目は \code{ERROR} で始まる行、2 行目は \code{txt} で終わる行、3 行目は行全体がちょうど \code{det\_A} である行だけを取る。特に 3 行目のように両端をアンカーで挟むと、\code{det\_AB} や \code{xdet\_A} が紛れ込むのを防げる。

\subsection{文字クラス \code{[...]} と代表的な省略記法}

\code{[abc]} は \code{a} か \code{b} か \code{c} の 1 文字、\code{[A-Z]} は英大文字 1 文字を表す。検出器名やセンサー番号のように、限られた候補から 1 文字だけ変わる文字列を扱うときに便利である。

\begin{lstlisting}[numbers=none, xleftmargin=2\zw]
^det_[AB]$
^run_[0-9][0-9][0-9]$
^[A-Za-z_][A-Za-z0-9_]*$
\end{lstlisting}

1 行目は \code{det\_A} または \code{det\_B}、2 行目は \code{run\_001} のような 3 桁番号付きの名前、3 行目は Python 風の識別子に近い形を表す。範囲指定の \code{[0-9]} は非常によく使う。

Python の \code{re} では、数字 1 文字を \code{\textbackslash d}、空白 1 文字を \code{\textbackslash s} と書ける。

\begin{lstlisting}[language=Python]
import re

re.fullmatch(r"det_[AB]", "det_A")
re.fullmatch(r"\d\d:\d\d:\d\d", "12:34:56")
re.search(r"\s+", "det_A   15.2")
\end{lstlisting}

ただし、\code{grep} や \code{awk} では \code{\textbackslash d} や \code{\textbackslash s} がそのまま使えないことがある。コマンドラインでは \code{[0-9]} や \code{[[:space:]]} を使う方が移植性が高い。

\subsection{量指定子 \code{*}, \code{+}, \code{?}}

量指定子は、直前の要素が何回くり返されるかを表す。\code{*} は 0 回以上、\code{+} は 1 回以上、\code{?} は 0 回または 1 回である。

\begin{lstlisting}[numbers=none, xleftmargin=2\zw]
ab*c
ab+c
colou?r
\end{lstlisting}

\code{ab*c} は \code{ac}、\code{abc}、\code{abbbc} に一致する。\code{ab+c} は少なくとも 1 個の \code{b} が必要なので \code{ac} には一致しない。\code{colou?r} は \code{color} と \code{colour} の両方に一致する。

解析では、「空白が 1 個以上続く」「符号が付くかもしれない」「小数点以下があるかもしれない」といった曖昧さを扱うときに量指定子が役立つ。

\begin{lstlisting}[language=Python]
pattern = re.compile(r"^[+-]?\d+(\.\d+)?$")

for text in ["12", "-3.5", "+0.25", "1.2.3"]:
    print(text, bool(pattern.fullmatch(text)))
\end{lstlisting}

このパターンは、符号付き整数や小数を表し、\code{1.2.3} のような壊れた文字列は弾く。

\section{物理・データ解析での具体例}

\subsection{ログから必要な行だけを探す}

実験ログやジョブログでは、まず異常行や特定の run だけを抜き出したい。アンカーと文字クラスを組み合わせると、かなり具体的な条件を書ける。

\begin{lstlisting}[numbers=none, xleftmargin=2\zw]
$ grep '^ERROR' run.log
$ grep '^2026-04-04 .* det_[AB] ' run.log
$ grep -E '^(WARN|ERROR)' run.log
\end{lstlisting}

3 行目では \code{grep -E} を使って \code{|} による選択を書いている。拡張正規表現を使うと括弧や \code{+} を素直に書けるので、条件が少し複雑になったらこちらの方が読みやすい。

\subsection{検出器 ID やセンサー名を選別する}

検出器名が \code{det\_A}, \code{det\_B}, \code{det\_C} のように並んでいるなら、文字クラスや選択を使って対象を絞れる。

\begin{lstlisting}[numbers=none, xleftmargin=2\zw]
$ grep '^det_[AB]$' detector_ids.txt
$ awk '$1 ~ /^det_[AB]$/ {print $1, $3}' summary.txt
\end{lstlisting}

\code{awk} では \code{\$1 ~ /.../} の形で、1 列目がその正規表現に一致するかどうかを判定できる。列を見ながら条件を書くときは \code{grep} より \code{awk} の方が自然なことが多い。

\subsection{Python で抽出と変換を同時に行う}

正規表現は「一致したかどうか」だけでなく、「どの部分が一致したか」を取り出すのにも向いている。Python では \code{re.search} や \code{re.fullmatch} の戻り値から、必要部分を取り出せる。

\begin{lstlisting}[language=Python]
import re

line = "2026-04-04 12:34:56 det_A signal=15.2"
match = re.search(r"signal=(\d+\.\d+)", line)

if match:
    signal = float(match.group(1))
    print(signal)
\end{lstlisting}

括弧 \code{(...)} で囲んだ部分はキャプチャされ、\code{group(1)}、\code{group(2)} のように取り出せる。ログ 1 行から run 番号、検出器名、信号値を抜き出して数値へ変換するような処理では、この形が便利である。

\section{Python の \code{re} と \code{grep}/\code{awk} の使い分け}

\subsection{Python では raw string を使う}

Python の文字列中で \code{\textbackslash} を多用すると読みにくくなる。そのため、正規表現を書くときは \code{r"..."} のような raw string を使うのが基本である。

\begin{lstlisting}[language=Python]
import re

pattern = re.compile(r"^det_[A-Z]\d+$")
print(bool(pattern.fullmatch("det_A12")))
\end{lstlisting}

raw string を使わないと、\texttt{"\textbackslash\textbackslash d+"} のようにバックスラッシュを二重に書かなければならない。正規表現の意図が見えにくくなるので、特別な事情がなければ raw string を選ぶ方がよい。

\subsection{\code{search}, \code{match}, \code{fullmatch} を使い分ける}

\code{re.search} は文字列のどこかに一致があればよいとき、\code{re.match} は先頭から一致していればよいとき、\code{re.fullmatch} は文字列全体がその形かを確かめたいときに使う。ファイル名や識別子を検証するときは、意図せず余分な文字が混ざらないよう \code{fullmatch} を使う方が安全である。

\begin{lstlisting}[language=Python]
import re

print(re.search(r"\d+", "run_42").group())
print(bool(re.match(r"run", "run_42")))
print(bool(re.fullmatch(r"run_\d+", "run_42")))
\end{lstlisting}

\subsection{コマンドラインで十分な場面と、Python に移すべき場面}

単に条件に合う行を見つけたい、件数を数えたい、特定列だけを抜きたい、といった用途なら \code{grep} や \code{awk} が速い。逆に、一致した部分を数値へ変換して次の計算へ渡したい、複数の条件分岐を組み合わせたい、抽出結果を DataFrame にしたい、といった場合は Python へ移した方が自然である。

\begin{itemize}
\item \code{grep}: 行全体の抽出や件数確認をすばやく行いたいとき
\item \code{awk}: 列条件と正規表現を組み合わせたいとき
\item Python \code{re}: 抽出した結果をそのまま解析処理へつなげたいとき
\end{itemize}

正規表現は強力だが、複雑な 1 行を無理に作ると自分でも読めなくなる。まずはアンカー、文字クラス、量指定子の 3 つで書けないかを考え、それでも足りないときに括弧や選択を足す方が安全である。解析コードと同じく、「まず読めること」を優先した方が長く使える。
